\section{CONVERGENCE OF THE HAUSDORFF SERIES (ULTRAMETRIC CASE)}

In this paragraph we assume that K is a non-discrete complete valued field of characteristic zero, with an ultrametric absolute value. We denote by p the characteristic of the residue field of K (Commutative Algebra, Chapter VI, § 3, no. 2).

If \(p \neq 0\) , we write \(a = |p|\) ; we know (Commutative Algebra, Chapter VI, §6, nos. 2 and 3) that \(0 < a < 1\) and that there exists one and only one valuation \(v\) on \(\mathbf{K}\) with values in \(\mathbf{R}\) whose restriction to \(\mathbf{Q}\) is the \(p\) -adic valuation \(v_p\) and which is such that \(|x| = a^{v(x)}\) for all \(x \in \mathbf{K}\) . Also we write:

\[
\theta = \frac {1}{p - 1}.\tag{1}
\]

If \(p = 0\) , we denote by \(a\) a real number such that \(0 < a < 1\) and by \(v\) a valuation on \(\mathbf{K}\) with values in \(\mathbf{R}\) such that \(|x| = a^{v(x)}\) for all \(x \in \mathbf{K}\) (loc. cit.). Then \(v(x) = 0\) for \(x \in \mathbf{Q}^*\) . Also we write:

\[
\theta = 0.\tag{2}
\]

